This work studied small lesion segmentation in brain MRI as a problem of
learning from sparse pathological signals. In multiple sclerosis MRI, small
lesions occupy a negligible fraction of the image volume, so a model
trained with a voxel-wise overlap objective can reach a high DSC while
leaving many of them undetected. The proposed CATMIL objective targets this
failure without any change to the network architecture: a
component-adaptive Tversky term gives small and large lesions comparable
weight, and a multiple-instance term penalizes any lesion in which no voxel
is detected.

CATMIL reduces the omission of small lesions. The improvement is
concentrated in the lesion sizes at which the standard losses fail, it is a
net gain in detected lesions rather than a trade between lesion sizes, and
it holds across test cases, across trained models, and on a second dataset
trained with the same loss weights. On MSLesSeg, CATMIL also preserves the
voxel-level overlap and boundary accuracy of the standard losses; on
3D-MR-MS its DSC is slightly lower. The analysis of predicted probabilities
shows why standard losses miss small lesions: they produce no response to
most of them, so no choice of decision threshold can recover these lesions.
The omission is therefore a consequence of how supervision is distributed
across lesions, rather than an inherent limitation of the network.

The cost of this sensitivity is a larger number of small false-positive
components, which lowers lesion-wise F1. These components are small and
localized, and a simple component-size filter removes most of them while
keeping the sensitivity gain. CATMIL is therefore most useful where missing
a lesion is a more serious error than a small false-positive component,
provided that the operating point is chosen after training for the
intended use.
